% !TEX root = ../main.tex
% Figure: Large-area coherent array synchronization mechanisms
% Chapter: ch20 - Large-Area Single-Mode Operation Strategies
% Description: Different coupling pathways for phase locking

\begin{figure}[htbp]
  \centering
  \begin{tikzpicture}[scale=0.92]
    % Array of emitters (simplified as circles)
    \def\emitterradius{0.3}
    
    % 5x5 emitter grid
    \foreach \x in {-4,-2,...,4} {
      \foreach \y in {-4,-2,...,4} {
        % Emitting aperture
        \fill[white] (\x,\y) circle (\emitterradius);
        \draw[thick, blue] (\x,\y) circle (\emitterradius);
        
        % Phase indicator (arrow showing oscillation phase)
        \pgfmathrandominteger{\phase}{0}{359}
        \draw[->, red, thick] (\x,\y) -- ++(\phase:0.2);
      }
    }
    
    % Background: Common substrate
    \fill[gray!10] (-5.5,-5.5) rectangle (5.5,5.5);
    \node[above right, gray] at (-5.5,-5.5) {公共衬底};
    
    % Coupling mechanism 1: Evanescent tunneling (nearest neighbors)
    \draw[very thick, green!60!black, <->] (-2,0) -- node[above] {\scriptsize 倏逝耦合} (0,0);
    \draw[very thick, green!60!black, <->] (0,0) -- (2,0);
    \draw[very thick, green!60!black, <->] (0,-2) -- (0,0);
    \draw[very thick, green!60!black, <->] (0,0) -- (0,2);
    
    % Coupling mechanism 2: Radiation coupling (long-range)
    \draw[very thick, orange, ->, bend left=20] (-4,-4) to[out=30,in=150] (4,4);
    \draw[very thick, orange, ->, bend right=20] (4,-4) to[out=150,in=30] (-4,4);
    \node[orange, align=center] at (0,3.5) {辐射耦合\\远场重叠};
    
    % Coupling mechanism 3: Shared cavity mode (global)
    \draw[ultra thick, purple, dashed] (0,0) circle (5);
    \node[purple, above] at (0,5.2) {共享腔模};
    
    % Boundary conditions
    \draw[double distance=2pt, thick] (-5.5,-5.5) rectangle (5.5,5.5);
    \node[below, font=\small] at (0,-5.8) {金属包覆边界或 DBR反射镜};
    
    % Zoom-in inset: Unit cell detail
    \begin{scope}[shift={(7,3)}, scale=2]
      % Single emitter with PhC cladding
      \fill[white] (0,0) circle (0.3);
      \draw[thick, blue] (0,0) circle (0.3);
      
      % Surrounding photonic crystal holes
      \foreach \angle in {0,60,...,300} {
        \fill[gray] (\angle:0.6) circle (0.08);
      }
      
      % Fundamental mode profile
      \draw[very thick, red, domain=0:360, samples=100] 
        plot({0.2*cos(\x)*(1 + 0.12*cos(2*\x))},{0.2*sin(\x)*(1 + 0.12*cos(2*\x))});
      
      \node[above, font=\small] at (0,0.8) {单胞放大};
      \node[below, align=center, font=\footnotesize] at (0,-0.8) {基模$\mathrm{LP}_{01}$\\近场高斯型};
    \end{scope}
    
    % Coherence area illustration
    \fill[yellow!20, opacity=0.5] (-1,-1) circle (1.5);
    \node[yellow!60!black, align=center] at (-1,-1) {相干区域\\$A_c$};
    
    % Key parameters box
    \node[draw, align=left, font=\small, anchor=north west] at (-6,6) {
      \textbf{同步条件}:\\
      $\bullet$ 耦合强度$\kappa >$频率失谐$\Delta\omega$\\
      $\bullet$ 相位锁定时间$\tau_c < $光子寿命$\tau_p$\\
      $\bullet$ 阵列尺寸$L <$相干长度$L_c$
    };
    
    % Output beam quality metrics
    \node[draw, align=left, font=\small, anchor=south east] at (6,-6) {
      \textbf{光束质量指标}:\\
      $\bullet$ $M^2 < 1.2$接近衍射极限\\
      $\bullet$ 旁瓣抑制比$>15\,\mathrm{dB}$\\
      $\bullet$ 指向稳定性$<0.1^\circ$
    };
  \end{tikzpicture}
  \caption{大面积 PCSEL 阵列中相干锁定的三种主要横向耦合机制：倏逝近邻耦合、辐射远场耦合以及共享全局腔模。右上角给出单胞局域模式示意，右下角列出评价大面积单模器件时常用的光束质量指标。}
  \label{fig:ch20_array_coupling}
\end{figure}
